\documentclass[11pt,a4paper]{article}
\usepackage{iftex}
\ifPDFTeX
  \usepackage[T1]{fontenc}
  \usepackage[utf8]{inputenc}
  \usepackage{lmodern}
  \input{glyphtounicode}
  \pdfgentounicode=1
\else
  \usepackage{fontspec}
  % Kerning off: kerned pairs ("A WS", "T ree") split words for ATS text extraction.
  \setmainfont{lmroman10-regular}[Extension=.otf, BoldFont=lmroman10-bold,
    ItalicFont=lmroman10-italic, BoldItalicFont=lmroman10-bolditalic, Kerning=Off]
\fi
\usepackage[margin=0.9in]{geometry}
\usepackage[hidelinks]{hyperref}
\hypersetup{pdftitle={Abhinav Kaushik - Cover Letter - Preply}, pdfauthor={Abhinav Kaushik}}

\pagestyle{empty}
\setlength{\parindent}{0pt}
\setlength{\parskip}{10pt}
\raggedright

\begin{document}

{\Large\bfseries Abhinav Kaushik}\\
Senior Applied AI Scientist\\
New Delhi, India \textbar{} +91 8700571859 \textbar{} \href{mailto:aabhinav.kaushik@gmail.com}{aabhinav.kaushik@gmail.com} \textbar{} \href{https://www.linkedin.com/in/abhinav-kaushik10/}{LinkedIn}  

September 27, 2026

Hiring Team\\
Preply

\textbf{Re: Senior Applied AI Scientist (AI Learning)}

Dear Hiring Team,

I am applying for the Senior Applied AI Scientist role on the AI Learning team. My work as a Senior Applied Scientist sits between research and production: fine-tuning transformer models, building RAG and conversational agents, and measuring whether they actually help. Learning has become the thread that ties that work to my life outside the office.

Preply's approach, where AI supports the tutor rather than replacing them, is the kind of educational AI I believe in, and personalised practice between lessons is a problem I have thought about for a long time. My fiancée lives in Berlin, and we plan to move to Barcelona together and settle there long term.

I would bring NLP work that has reached real users. I fine-tuned a cross-encoder to re-rank dense retrieval over survey data, and built a multi-turn conversational chatbot in LangGraph that cut manual analysis by 98\%. I also designed LLM evaluators and tracing that reduced human review by 43\%, which is how I would approach measuring learner progress.

Outside work, I am building ThinkTree.AI, a non-profit AI learning platform used by NGO teachers to support students with ADHD. Watching how those students learn has shaped how I think about personalisation and patience in AI tools. I would love to talk about what I could bring to your team.

Sincerely,\\[4pt]
Abhinav Kaushik

\end{document}